\documentclass[11pt]{article}
\usepackage[margin=0.9in]{geometry}
\usepackage[colorlinks=true,urlcolor=blue]{hyperref}
\pagestyle{empty}
\begin{document}

\begin{flushleft}
Quang-Vinh Dang\\
British University Vietnam\\
Hung Yen, Vietnam\\
\href{mailto:vinh.dq4@buv.edu.vn}{vinh.dq4@buv.edu.vn}\\[12pt]
\today\\[12pt]
\end{flushleft}

\noindent To the Editors, \emph{Journal of Engineering and Technology Management}\\
Special Issue: ``Evolutionary Perspective of the Dynamics of Industries: Revisiting the Role of Social Technologies''

\vspace{10pt}
\noindent Dear Editors,

\vspace{6pt}
\noindent Please find enclosed our manuscript, \textbf{``Anytime-valid monitoring of technology adoption dynamics: Digital technology diffusion surveillance in the United States, 1990--2026,''} for consideration for the special issue.

The paper addresses a problem that has become acute as technology-adoption data have shifted from annual surveys to continuously updated streams: dashboards and analysts now inspect adoption series at every release, but the classical statistical toolkit of diffusion research assumes a single test on a closed sample, and its guarantees dissolve under repeated inspection. We develop an anytime-valid surveillance framework that (i) treats survey-instrument redesigns as causal interventions on the measurement process, distinct from genuine shifts in the underlying diffusion state, and (ii) uses one-sided mixture e-processes whose error control holds uniformly over time, so that continuous monitoring and data-dependent stopping do not invalidate the inference.

We apply the framework to the U.S. Census Bureau's Business Trends and Outlook Survey, the highest-frequency official adoption series ever published for a general-purpose technology, tracking enterprise AI use biweekly since September 2023. The monitor declares a diffusion takeoff on 20 April 2025, quantifies the widely discussed mid-2025 growth deceleration as strong but sub-threshold evidence that is subsequently vindicated by re-acceleration, and correctly absorbs a 25.7-standard-error question-wording redesign in November 2025 as a measurement event rather than a diffusion shock -- an event that conventional repeated-testing rules misattribute as the largest ``adoption shock'' on record. We further extend the surveillance logic (with an appropriately caveated plug-in-variance construction) across three decades of U.S. digital diffusion, recovering the internet's 1995 takeoff, 2008 saturation onset, and e-commerce's 2004, 2018, 2020, and 2023 regime shifts as dated declarations, and place the intervening technology waves (data mining, cloud, big data, data science, deep learning) on the same timeline using canonical landmarks, since no official adoption series exist for them.

We believe the paper fits the special issue's call directly: it operationalizes an evolutionary view of technology dynamics by tracking punctuated diffusion as it unfolds, as a sequence of dated, error-controlled regime declarations, rather than fitting a curve to data after the fact.

All data used are public (U.S. Census Bureau BTOS statistics and FRED series); the exact analysis datasets and replication scripts reproducing every reported number, table, and figure are included as supplementary material. The manuscript is original, has not been published elsewhere, and is not under consideration by any other journal. We have no conflicts of interest to disclose.

Thank you for considering our manuscript. We look forward to your response.

\vspace{14pt}
\noindent Sincerely,\\[6pt]
Quang-Vinh Dang\\
British University Vietnam\\
Corresponding author

\end{document}
